\section{Reliability and bounded EXIT}
\label{sec:reliability}

Let $\cH$ be a repair clutter on helper set $V$.  For $v\in V$, define
\[
 \cH-v=\min\{E\in\cH:v\notin E\},
 \qquad
 \cH/v=\min\{E\setminus\{v\}:E\in\cH\}.
\]

\begin{theorem}[Reliability calculus]
\label{thm:reliability}
For independent helper survival probabilities $(s_v)$,
\[
 R_\cH(s)=(1-s_v)R_{\cH-v}(s)+s_vR_{\cH/v}(s),
\]
and
\[
 \frac{\partial R_\cH}{\partial s_v}
 =R_{\cH/v}-R_{\cH-v}
 =\Pr(v\text{ is pivotal}).
\]
Under homogeneous survival $s$,
$R_\cH'(s)=\sum_v\Pr_s(v\text{ is pivotal})$.
If the minimum blocker size is $\tau$ and there are $b_\tau$ minimum
blockers, then for helper failure probability $p$,
\[
 1-R_\cH(1-p)=b_\tau p^\tau+O(p^{\tau+1}).
\]
\end{theorem}

\begin{proof}
Write \(f_\cH(A)=1\) when \(A\) contains an edge of \(\cH\), and write it as
zero otherwise.  The reliability is the finite multilinear sum
\[
 R_\cH(s)=\sum_{A\subseteq V}f_\cH(A)
 \prod_{y\in A}s_y\prod_{y\notin A}(1-s_y).
\]
Partition the subsets according to whether they contain \(v\).  Erasing \(v\)
from those that do contain it identifies the two parts with the Boolean
systems \(\cH-v\) and \(\cH/v\), respectively.  This gives
\[
 R_\cH=(1-s_v)R_{\cH-v}+s_vR_{\cH/v}.
\]
The two conditional reliabilities do not involve \(s_v\), so differentiation
gives \(R_{\cH/v}-R_{\cH-v}\).  Monotonicity makes this difference the
conditional probability that adjoining \(v\) changes failure into success,
which is precisely pivotality.  On the diagonal \(s_v=s\), the chain rule
sums these coordinate derivatives and gives the Russo--Margulis identity.

Now expose the failed set \(F\).  Repair fails exactly when \(F\) is a
transversal of \(\cH\), equivalently when it contains a minimal blocker.
Thus
\[
 1-R_\cH(1-p)=
 \sum_{\substack{F\subseteq V\\F\ {\rm is\ a\ transversal}}}
 p^{|F|}(1-p)^{|V|-|F|}.
\]
There are no summands below degree \(\tau\), and the transversals of size
\(\tau\) are exactly the \(b_\tau\) minimum blockers.  Their contribution is
\(b_\tau p^\tau(1-p)^{|V|-\tau}
=b_\tau p^\tau+O(p^{\tau+1})\).  Every remaining term contains at least
\(p^{\tau+1}\); the sum is finite, proving the expansion.
\end{proof}

These are the standard deletion--contraction and pivotality identities for a
monotone reliability system \cite{Colbourn1987}.

For EXIT conventions, condition the target $x$ unavailable and erase helper
$v$ with probability $p_v$.  Put
\[
 h_x^{(\leq r)}(p)=\Pr(f_x^{(\leq r)}(A)=0),
 \qquad R_x^{(\leq r)}=1-h_x^{(\leq r)}.
\]
This is an \emph{extrinsic} failure probability.  If the target is itself
erased with probability $p_x$, its residual erasure probability is
$p_xh_x^{(\leq r)}$.  At full radius it is symbol-MAP EXIT; at finite radius
it is a bounded-query decoder and carries no capacity claim
\cite{AshikhminKramerTenBrink2004,Mazumdar2023}.

\begin{proposition}[Bounded-EXIT hierarchy]
\label{prop:bounded-exit}
The failure form of deletion--contraction is
\[
 h_\cH(p)=p_vh_{\cH-v}(p)+(1-p_v)h_{\cH/v}(p),
 \qquad
 \frac{\partial h_\cH}{\partial p_v}
 =\Pr(v\text{ is pivotal}).
\]
If $L_x(A)$ is the size of the cheapest repair contained in the realized
survivor set, with $L_x=\infty$ when none exists, then
\[
 \Pr(L_x=r)=h_x^{(\leq r-1)}-h_x^{(\leq r)},
 \qquad
 \Pr(L_x=\infty)=h_x^{\rm MAP}.
\]
The first identity is for \(r\geq1\).
\end{proposition}

\begin{proof}
Conditioning on whether \(v\) is erased gives the displayed recurrence:
when it is erased, only repairs avoiding \(v\) remain; when it survives, it
may be contracted from every repair.  Equivalently, substitute
\(s_y=1-p_y\) in the reliability recurrence and take complements.
Differentiation reverses both signs, so the derivative is again the pivotal
probability.

Let \(\cH_x^{(\leq r)}\) be the full port filtered by edge size at most \(r\).
The event \(L_x\leq r\) is exactly success for this truncated port.  Hence
\[
 \{L_x=r\}=\{L_x>r-1\}\setminus\{L_x>r\},
\]
and taking probabilities gives the difference of failure curves.  At full
radius, the complementary event is that no repair exists, which is
\(L_x=\infty\).
\end{proof}

Minimal transversals of $\cH_x^{(\leq r)}$ are exact target-specific,
one-shot bounded-repair failure certificates.  They are not automatically
the representation-dependent stopping sets of a global iterative Tanner
decoder.

For a length-$N$ code define the radius-$r$ locality deficit
\[
 L_r(x)=\int_0^1\bigl(h_x^{(\leq r)}(p)-h_x^{\rm MAP}(p)\bigr)\,dp.
\]
If $\Delta a_k$ counts size-$k$ survivor sets repaired by full MAP but not at
radius $r$, beta integration gives
\[
 L_r(x)=\sum_{k=0}^{N-1}
 \frac{\Delta a_k}{N\binom{N-1}{k}}.
\]
With entropy normalized in $\log_q$ units, the unrecoverable-symbol indicator
is the normalized conditional entropy on the $q$-ary erasure channel.  Under
this convention the symbol-MAP areas sum to the code dimension $K$
\cite{AshikhminKramerTenBrink2004}, so
\[
 \sum_x\int_0^1h_x^{(\leq r)}(p)\,dp
 =K+\sum_xL_r(x).
\]
